% ---
% title: Nabla Operator
% date: 2026-04-01
% tag: Notes
% excerpt: 哈密顿算子的三个用法：梯度、散度和旋度，配上直角坐标下的定义和一点几何直觉。
% ---
设 \(\hat{x},\hat{y},\hat{z}\) 为三维空间内的一组正交基，哈密顿算子定义为：

\[
\nabla=\hat{x}\frac{\partial}{\partial x} + \hat{y}\frac{\partial}{\partial y} + \hat{z}\frac{\partial}{\partial z}
\]

它含有微分算子，因此可以描述数学和物理对象在空间中的瞬时变化，例如函数值的变化或流场的局部变化。也可以把它理解为欧氏空间 \((\hat{x},\hat{y},\hat{z})\) 中，描述对象沿某一方向的瞬时变化率。

下面是我对梯度、散度和旋度的直观理解。

\section{梯度}

梯度描述标量函数 \(T=t(x,y,z)\) 在某一点处函数值上升最快的方向。它是一个矢量，定义为：

\begin{align}
\nabla T 
&=\hat{x}\frac{\partial t(x,y,z)}{\partial x} + \hat{y}\frac{\partial t(x,y,z)}{\partial y} + \hat{z}\frac{\partial t(x,y,z)}{\partial z}
\end{align}

也就是说，梯度给出函数在欧氏空间中的变化趋势，方向指向函数值增大的一侧。

\section{散度}

散度是哈密顿算子与一个矢量场做点乘，结果是一个标量。这个标量反映各分量中彼此相关的部分的累积。类比做功：力与位移做内积时，垂直方向的分量不计入，只有相关的分量相乘后累加。

\begin{align}
\nabla \cdot v &=\big(\hat{x}\frac{\partial}{\partial x} + \hat{y}\frac{\partial}{\partial y} + \hat{z}\frac{\partial}{\partial z}\big)\cdot (v_x\hat{x}+v_y\hat{y}+v_z\hat{z}) \\
&= \frac{\partial v_x}{\partial x}+ \frac{\partial v_y}{\partial y} + \frac{\partial v_z}{\partial z}\\
&=空间 \cdot 向量场
\end{align}

散度是一个矢量场（比如流场）在空间各方向上累积后得到的净值。对流体来说，散度大于 0 表示该点附近有净流出，小于 0 表示净流入。

\section{旋度}

旋度也可以从流体的角度理解。它是哈密顿算子与矢量场做叉乘，结果仍是矢量，定义为：

\begin{align}
\nabla \times \mathbf{v}
&=\left(\hat{x}\frac{\partial}{\partial x} + \hat{y}\frac{\partial}{\partial y} + \hat{z}\frac{\partial}{\partial z}\right)\times
\left(v_x\hat{x}+v_y\hat{y}+v_z\hat{z}\right) \\
&=
\begin{vmatrix}
\hat{x} & \hat{y} & \hat{z} \\
\dfrac{\partial}{\partial x} & \dfrac{\partial}{\partial y} & \dfrac{\partial}{\partial z} \\
v_x & v_y & v_z
\end{vmatrix} \\
&=空间 \times 向量场
\end{align}

哈密顿算子作用在具体函数上，这里的空间不是物理空间，而是函数所在的抽象空间。叉乘与点乘的差别在于，它描述的是矢量场在法向上的累积（点乘对应切向），同样可以用流体类比。如果旋度指向某个方向，按右手定则，流场会绕该方向旋转，旋转的中心就是涡旋的中心。
